\documentclass[11pt]{article}
\usepackage{setspace}
\usepackage[top=1in, bottom=1in, left=1in, right=1in]{geometry}
\usepackage{hyperref}
\onehalfspacing

% Title information
\title{Preferences on Tax Enforcement and Tax Administration:\\Note on Information Treatments}
\author{Tim Kircali, Stefanie Stantcheva, Matthias Weber}
\date{\today} % Use \date{} for no date

% Bibliography package
\usepackage[backend=biber, style=apa]{biblatex}
\addbibresource{references.bib} % Link to the bibliography file

\begin{document}
\maketitle
\section{Information Treatments}
\subsection{Tax Evasion Treatment}
\subsubsection*{General}
The first treatment provides respondents with information about the economic effects of third-party reporting and AEoI in tax enforcement. We give them information about the size of the tax gap and we directly point out that third-party reporting and AEoI is effective in reducing tax evasion by explaining evidence which shows that tax evasion is rare in areas which are covered under third-party reporting and AEoI. The primary goal of the tax evasion treatment is to manipulate people’s knowledge about the economic effects of tax enforcement, third-party reporting and AEoI without affecting considerations related to fairness \& inequality.   

\subsubsection*{Information Text}
The following explanations could be used as subtitles and audio for an animated video.
\begin{itemize}
\item One problem governments face is the loss of tax revenue due to tax evasion. This means that some taxpayers deliberately hide some of their income or wealth from the tax administration. 
\item In the US, for example, the Internal Revenue Service estimates that only 85\% of taxes are paid voluntarily and on time. The remaining 15\% amount to \$696 billion and more than \$600 billion of these are never collected
\begin{itemize}
    \item Animation: These numbers will be visualized with diagrams.
\end{itemize}
\item The tax authorities have a number of options to combat tax evasion.
\item For example, they can increase their audit activities, which means that they investigate more tax returns or taxpayers in detail.  
\item Research has shown that tax evasion hardly takes place in areas that are covered by third-party reporting.
\item Third-party reporting means the tax authority uses information from other sources --- like employers, banks or insurance companies. This information is automatically sent from domestic firms and institutions to the tax authority.
\item For example, employers report to the tax administration how much they pay their employees in wages.
\item Another example is that a bank automatically informs the tax authority about interest and dividend payments of assets held by a person in a bank.
\item Increasing the amount of data that is third-party reported is therefore an effective way to reduce tax evasion.
\begin{itemize}
    \item Animation: We can visualize two parts of financial data, one that is third-party reported and one that is not. We show that most of the tax gap is the result of the unreported part. We then show that the tax gap decreases as we increase the proportion that is third-party reported.
\end{itemize}
\item Besides domestic third-party reporting, international information exchange is also very effective in reducing tax evasion. 
\item Several countries started to automatically exchange information on bank accounts held by foreign residents. This has helped tax authorities to detect taxpayers who are hiding income and wealth abroad with the intention to evade taxes.
\item This automatic exchange of information has proven to be very effective. For example, in Denmark, the automatic exchange of information has led to an estimated reduction of offshore tax evasion of about 70\%. 

\end{itemize}
\subsubsection*{Follow up questions}
We will ask 2-3 very basic questions to understand if respondents actually watched the video.
\subsection{Inequality Treatment}
\subsubsection*{General}
The second specific information treatment again includes information on tax evasion, but also explains the fact that the tax gap is much more pronounced for high-income individuals and discusses why this is the case. After that, we discuss fairness considerations of tax evasion, such as the fact that honest taxpayers have to finance the fiscal gap caused by tax evaders. The primary goal of the inequality treatment is to manipulate people’s knowledge about distributional consequences of third-party reporting and AEI.  
\subsubsection*{Information Text}
The following explanations could be used as subtitles and audio for an animated video.
\begin{itemize}
    \item One problem governments face is the loss of tax revenue due to tax evasion. This means that some taxpayers deliberately hide some of their income or wealth from the tax administration. 
    \item In the US, for example, the Internal Revenue Service estimates that only 85\% of taxes are paid voluntarily and on time. The remaining 15\% amount to \$696 billion and more than \$600 billion of these are never collected
    \begin{itemize}
        \item Animation: These numbers will be visualized with diagrams.
    \end{itemize}
    \item Research has shown that not all groups of the population evade taxes to the same extent. 
    \item First, taxpayers with low to medium incomes generally only evade a small share of their income.
    \begin{itemize}
        \item Animation: With a bar chart for 10 income groups, where each bar has the same height, we can visualize evasion rates.  
        \item Here we could use examples of \cite{johns2010distribution}, which show evasion rates along the income distribution for deciles and quantiles.
    \end{itemize}
    \item Second, since wealthy taxpayers generally have to pay more taxes, the evaded share of their income weighs heavily. 
    \begin{itemize}
        \item Animation: Using the same bar chart we can visualize by growing the bars that higher income groups' tax evasion is much higher in monetary terms.
    \end{itemize}
    \item Therefore, the large amount of total evaded taxes is mainly the result of tax evasion among wealthy taxpayers.
    \item That means a small group of the population is responsible for a large part of missing taxes. For example (x)\% of the absolute tax gap is due to the (y)\% richest part of the population. 
    \item Recent research has shown that tax evasion is even unequally distributed among these mentioned (y)\%. 
    \begin{itemize}
        \item Animation: We zoom in to the bar of the (y)\%. (It is important for respondents to be able to keep track of where along the entire income distribution the information is shown.) 
    \end{itemize}
    \item An analysis of leaks of undeclared offshore bank accounts shows that that only people at the very top of the income distribution are using offshore bank accounts to evade taxes. 
    \begin{itemize}
        \item Animation: We show that tax evasion is also highly unequally distributed at the top of the income distribution, with most of the evasion driven by the top 1\%.
        \item For that we can show evidence by \cite{alstadsaeter2019tax} where evasion rates are even higher than suggested in \cite{johns2010distribution}. 
    \end{itemize} 
\item The researchers also estimate that the super-rich, referring to the top 0.01 per cent of Scandinavian taxpayers, evade about 25 per cent of their tax liability just by hiding assets and investment income abroad. 
\begin{itemize}
    \item Animation: We show previously suggested evasion rates by quantile and show that the evasion estimate for top income by \cite{alstadsaeter2019tax} suggests even higher evasion rates. 
    \item We might also consider not showing this information if previous info is already convincing enough. 
\end{itemize}
\item When a group of taxpayers evades taxes, the government has less money to spend on public services. It must either cut services or raise taxes. This will harm honest taxpayers, who have paid their fair share of taxes. 
\begin{itemize}
    \item Animation: Show the government budget and the government budget gap. Animate how that affects honest taxpayers.
    \item We might also consider not showing this information if previous info is already convincing enough. 
\end{itemize}
\item The tax authorities have a number of options to combat tax evasion.
\item For example, they can increase their audit activities, which means that they investigate more tax returns or taxpayers in detail.  
\item Research has shown that tax evasion hardly takes place in areas that are covered by third-party reporting.
\item Third-party reporting means the tax authority uses information from other sources --- like employers, banks or insurance companies. This information is automatically sent from domestic firms and institutions to the tax authority.
\item For example, employers report to the tax administration how much they pay their employees in wages.
\item Another example is that a bank automatically informs the tax authority about interest and dividend payments of assets held by a person in a bank.
\item Increasing the amount of data that is third-party reported is therefore an effective way to reduce tax evasion.
\begin{itemize}
    \item Animation: We can visualise two parts of the financial data, one that is third-party reported and one that is not. We show that most of the tax gap is the result of the unreported part. We then show that the tax gap decreases as we increase the proportion that is third-party reported.
\end{itemize}
\item Besides domestic third-party reporting, international information exchange is also very effective in reducing tax evasion.  
\item Several countries started to automatically exchange information on bank accounts held by foreign residents. This has helped tax authorities to detect taxpayers who are hiding income and wealth abroad with the intention to evade taxes.
\item This automatic exchange of information has proven to be very effective. For example, in Denmark, the automatic exchange of information has led to an estimated reduction of offshore tax evasion of about 70\%.
\end{itemize}

%\subsection{Data Privacy Treatment}
%\subsubsection*{General}
%In the data privacy treatment, we explain that the IRS already has much tax-relevant data and that taxpayers only had to give up a small last share data privacy to help them to fulfill its functions. The primary goal of the data privacy treatment is to manipulate people’s knowledge about third-party reporting and AEoI, with focus on data privacy considerations.
%
%\subsubsection*{Information Text}
%The following explanations could be used as subtitles and audio for an animated video.
%\begin{itemize}
%    \item Governmental institutions collect data on citizens in order to fulfill their functions and provide their services.
%    \item For instance, the Internal Revenue Service does not rely solely on taxpayers' self-reported information. It also collects information reported by third parties to verify the correctness of the self-reported information.
%    \item Third-party reporting means the tax authority uses information from other sources --- like employers, banks or insurance companies. This information is automatically sent from domestic firms and institutions to the tax authority.
%    \item For example, employers report to the tax administration how much they pay their employees in wages.
%    \item Furthermore, banks automatically inform the tax authority about interest and dividend payments of assets held by a person in a bank.
%    \item But the IRS does not automatically obtain all tax-relevant data it could use. 
%    \item Often, other governmental institutions hold tax-related information but these data are not shared with the IRS.
%    \item For example, millions of people receive health insurance coverage through Medicaid or other governmental programs. But this information often stays within the health agencies and is not shared with the IRS. Not receiving these data makes it harder for the IRS to make sure that people receive the correct amounts of tax credits.
%    \item Another example is that local governments keep track of who owns a house and how much property tax they pay. But they do not send this information to the IRS. This means the IRS does not automatically know if someone owns a home or rents it out, which would be relevant for the IRS to assess housing-related tax deductions.
%    %\item But this is tax-relevant information which is important for the IRS. 
%    \item For these data, the IRS needs to rely on taxpayers' declarations or to conduct tax audits verifying the information. This increases the likelihood of errors and leads to administrative costs and burden.
    %\begin{itemize}
    %\item Unemployment benefits: tax relevant data actually should be sent to the IRS but often does not happen because of administrative problems.
    %\item Example text for unemployment benefits: When someone gets unemployment benefits, that information stays with the state. The IRS is supposed to get a copy, but sometimes states are late, send the wrong data, or don’t send it at all. That means the IRS may not know a person received unemployment income — unless they report it themselves.
    %\item Millions of people get free or low-cost health insurance through Medicaid or the government health insurance marketplace. But that information stays with state health agencies or healthcare.gov. The IRS doesn’t always receive it — or gets it late or with errors. That makes it harder to make sure people get the right tax credits, or to verify who had coverage.
    %\item Medicaid is not taxable, but the IRS could still use that information. For example, it could help check whether someone qualifies for certain tax credits, or avoid giving the wrong kind of subsidy. It could also help the IRS detect mistakes or fraud, and make the tax filing process easier for low-income people
    %\item Local governments keep track of who owns a house and how much property tax they pay. But they don’t send this information to the IRS. That means the IRS doesn’t automatically know if someone owns a home, rents it out, or claims tax deductions related to it — unless the person reports it themselves.
    %\end{itemize}
    %\item Thanks to the existing third-party reporting, the IRS already has most of the tax-relevant data of taxpayers.
    %\begin{itemize}
    %    \item Animation: We visualize the amount of data the IRS already has with a colour (e.g. green).
    %\end{itemize}
    %\item However, the IRS has no direct access to some tax-relevant data. For this share of the data, the tax administration must either rely on taxpayers' self-reported data or conduct tax audits to verify information.
    %\begin{itemize}
    %    \item Animation: We visualize the amount of data the IRS has no direct access with another color (e.g. red).
    %\end{itemize}
%    \item It may be more efficient if all tax-relevant information were consistently passed on to the IRS. This would help the IRS to fulfill its functions and provide services to taxpayers.
%    \item (For example, the tax authorities can better fight tax evasion if they have all tax relevant information.) 
%   \item (The tax authority would also be able to prepopulate tax returns if they have all necessary data. Then taxpayers would not need to fill in the tax returns themselves but only check and confirm the tax returns prepopulated by the tax authority.) 
    %\item Most of the taxpayers would not even have to worry about privacy, since most of the relevant data is already reported to the IRS.
    
    %\item If taxpayers were willing to share all tax-relevant information with the authorities, this would considerably help them fulfil their functions and provide their services. 
    %\item Most taxpayers would not even have to give up much privacy, since most of the relevant data is already reported to the IRS.
%\end{itemize}
%\subsubsection*{Follow up questions}
%We will ask 2-3 very basic questions to understand if respondents actually watched the video.
%\subsection{Tax Filing Treatment}
%\subsubsection*{General}
%The last treatment discusses concerns about tax filing costs such as the fact that an average American taxpayer has a tax filing effort and cost of about 13 hours + 270 \$. It describes how (guided) tax software works and what a prepopulated tax return looks like. We emphasize that tax software and prepopulated tax returns help tax filers to make fewer mistakes and significantly reduce their tax filing costs. Importantly, we highlight the importance of third-party reporting to the tax administration in providing prepopulated tax returns. The primary goal of the tax filing treatment is to manipulate people’s knowledge and perceptions about tax filing costs and tax filing services and how they are related to the existence of third-party reporting and AEI.  
%
%\subsubsection*{Information Text}
%The following explanations could be used as subtitles and audio for an animated video.
%\begin{itemize}
%    \item Most people need to file tax returns to report information to the tax administration, such as received salary, bank account balances or paid contributions to health insurance.
%    \item While some taxpayers prepare their tax returns themselves, others pay for advice or pay someone else to prepare their tax returns. 
%    \item Taxpayers who prepare their tax returns themselves often use specific tax preparation software to file their taxes. Such software helps taxpayers file their tax returns by providing guidance through the tax formulas and explaining which deductions might apply. It provides field-by-field advice to reduce errors and foregone reductions.
%    \begin{itemize}
%        \item Animation: We will visualize how tax software works, such that people get an understanding of it. 
%    \end{itemize}
%    \item All these options are time consuming or costly. An estimate by the US tax administration, the Internal Revenue Service, suggests that taxpayers in the US need on average  13 hours plus \$270 for tax software and help to file their taxes.
%    \item While most taxpayers in the US have to pay for tax software licenses, many other countries offer free government-provided tax software, which makes it easier for taxpayers to file their taxes. This would also be possible for the US. 
%    \begin{itemize}
%        \item Animation: We make it visually clear that tax payers in some countries are provided with better resources, while others are not. 
%    \end{itemize}
%    \item The IRS could potentially also offer prepopulated tax returns. This means that the tax returns are pre-filled with information the tax administration already has. Prepopulated tax returns make the process of filing taxes much shorter and easier, with lower potential for making mistakes.
%    \begin{itemize}
%        \item Animation: We show what a prepopulated tax return looks like and clearly visualise that many/most fields are already prepopulated, thus reducing the taxpayer's filing burden. 
%    \end{itemize}
%    \item If the tax administration does not have the respective data, it will not be able to prepopulate many fields, which will then have to be filled in manually by taxpayers. The more tax-relevant information the tax administration has, the more fields in the tax return can be prepopulated.
%    \begin{itemize}
%        \item Animation: We clearly illustrate that many fields remain empty when the level of information is low, and that taxpayers then have to file manually.
%    \end{itemize}
%    \item The tax authority often gets access to tax-relevant information by third-party reporting. 
%    \item Third-party reporting means the tax authority uses information from other sources - like employers, banks or insurance companies. This information is automatically sent from domestic firms and institutions to the tax authority.
%    \item The more information the tax administration automatically receives, the better it can provide pre-populating services.
%    \begin{itemize}
%        \item Animation: We can show visually how more data translates into more fields being pre-filled. 
%    \end{itemize}
%\end{itemize}
%\subsubsection*{Follow up questions}
%We will ask 2-3 very basic questions to understand if respondents actually watched the video.
%\subsection{Full Info Treatment}
%\subsubsection*{General}
%The full information treatment includes information from all treatments. It can be interpreted as a balanced video including the highlights of all the treatments. It helps us to understand how information transmission is changing if the information is presented next to each other. 
%\subsubsection*{Information Text}
%\begin{itemize}
%    \item One problem governments face is the loss of tax revenue due to tax evasion. This means that some taxpayers deliberately hide some of their income or wealth from the tax administration. 
%    \item In the US, for example, the Internal Revenue Service, estimates that only 85\% of taxes are paid voluntarily and on time. The remaining 15\% amount to \$696 billion and more than \$600 billion of these are never collected
%    \item Research has shown that this big amount of evaded taxes is mainly result of tax evasion among wealthy taxpayers.
%    \item That means a small group of the population is responsible for a large part of the tax gap. For example (x)\% of the absolute tax gap is due to the (y)\% richest part of the population. 
%    \item The tax authorities have a number of options to combat tax evasion.
%    \item For example, they can increase their audit activities, which means that they investigate more tax returns or taxpayers in detail.  
%    \item Research has shown that tax evasion hardly takes place in areas that are covered by third-party reporting.
%    \item Third-party reporting means the tax authority uses information from other sources - like employers, banks or insurance companies. This information is automatically sent from domestic firms and institutions to the tax authority.
%    \item For example, employers report to the tax administration how much they pay their employees in wages.
%    \item Furthermore, banks automatically inform the tax authority about interest and dividend payments of assets held by a person in a bank. 
%    \item Increasing the amount of data that is third-party reported is therefore an effective way to reduce tax evasion.
%    \item Besides domestic third-party reporting, international information exchange is also very effective in reducing tax evasion. 
%    \item Tax authorities of several countries started to automatically exchange information of foreign bank accounts. This has helped them to detect taxpayers who are hiding income and wealth abroad with the intention to evade taxes.
%    \item Thanks to the existing third-party reporting, the IRS already has most of the tax-relevant data of taxpayers.
%    \item However, the IRS has no direct access to some tax-relevant data. For this share of the data, the tax administration must either rely on taxpayers' self-reported data or conduct tax audits to verify information.
%    \item Access to this final portion of data would be extremely valuable for the tax administration, as some functions and services require the complete set of information.
%     \item If taxpayers were willing to share all tax-relevant information with the authorities, this would considerably help them fulfil their functions and provide their services. 
%    \item Most taxpayers would not even have to give up much privacy, since most of the relevant data is already reported to the IRS.
%    \item And taxpayers might also gain directly because the tax administration would be able to reduce tax filing costs. An estimate by the US tax administration, the Internal Revenue Service, suggests that taxpayers in the US need on average 13 hours plus \$270 for tax software and help to file their taxes.
%    \item One way to reduce these costs is to offer free tax software. 
%    \item The IRS could also potentially go a step further and offer prepopulated tax returns. This means that the tax returns are pre-filled with information the tax administration already has. Prepopulated tax returns make the process of filing taxes much shorter and easier, with lower potential for making mistakes.
%    \item However, if the tax administration does not have the respective data, it will not be able to prepopulate much, and most fields will be left blank and have to be filled in manually by taxpayers. 
%\end{itemize}
%\subsubsection*{Follow up questions}
%We will ask 2-3 very basic questions to understand if respondents actually watched all parts of the video.
%

\subsection{Data and Filing Treatment}
\subsubsection*{General}
The third treatment discusses concerns about tax filing costs such as the fact that an average American taxpayer has a tax filing effort and cost of about 13 hours + 270 \$. It describes how (guided) tax software works and what a prepopulated tax return looks like. We emphasize that tax software and prepopulated tax returns help tax filers to make fewer mistakes and significantly reduce their tax filing costs. Importantly, we highlight the importance of third-party reporting to the tax administration in providing prepopulated tax returns. We explain that the IRS already has much tax-relevant data and that taxpayers only had to give up a small last share data privacy to help them to fulfill its functions. 
\subsubsection*{Information Text}
The following explanations could be used as subtitles and audio for an animated video.
\begin{itemize}
    \item Most people need to file tax returns to report information to the tax administration, such as received salary, bank account balances or paid contributions to health insurance.
    \item While some taxpayers prepare their tax returns themselves, others pay for advice or pay someone else to prepare their tax returns. 
    \item Taxpayers who prepare their tax returns themselves often use specific tax preparation software to file their taxes. Such software helps taxpayers file their tax returns by providing guidance through the tax formulas and explaining which deductions might apply. It provides field-by-field advice to reduce errors and foregone reductions.
    \begin{itemize}
        \item Animation: We will visualize how tax software works, such that people get an understanding of it. 
    \end{itemize}
    %\item All these options are time consuming or costly. An estimate by the US tax administration, the Internal Revenue Service, suggests that taxpayers in the US need on average  13 hours plus \$270 for tax software and help to file their taxes.
    \item While most taxpayers in the US have to pay for tax software licenses, many other countries offer free government-provided tax software, which makes it easier for taxpayers to file their taxes. This would also be possible for the US. 
    \begin{itemize}
        \item Animation: We make it visually clear that tax payers in some countries are provided with better resources, while others are not. 
    \end{itemize}
    \item The Internal Revenue Service, short the IRS, could potentially also offer prepopulated tax returns. This means that the tax returns are pre-filled with information the tax administration already has. Prepopulated tax returns make the process of filing taxes much shorter and easier, with lower potential for making mistakes.
    \begin{itemize}
        \item Animation: We show what a prepopulated tax return looks like and clearly visualise that many/most fields are already prepopulated, thus reducing the taxpayer's filing burden. 
    \end{itemize}
    \item In order to offer pre-populated tax returns, the IRS needs the relevant tax data. Without this data, it will not be able to prepopulate many fields. Taxpayers will then have to fill in these fields manually.
    %\item If the tax administration does not have the respective data, it will not be able to prepopulate many fields, which will then have to be filled in manually by taxpayers. The more tax-relevant information the tax administration has, the more fields in the tax return can be prepopulated.
    \begin{itemize}
        \item Animation: We clearly illustrate that many fields remain empty when the level of information is low, and that taxpayers then have to file manually.
    \end{itemize}
    \item The tax authority often gets access to tax-relevant information by third-party reporting. 
    \item Third-party reporting means the tax authority uses information from other sources - like employers, banks or insurance companies. This information is automatically sent from domestic firms and institutions to the tax authority.
    \item For example, employers report to the tax administration how much they pay their employees in wages.
    \item Furthermore, banks automatically inform the tax authority about interest and dividend payments of assets held by a person in a bank.
    \item But the IRS does not automatically obtain all tax-relevant data it could use. 
    \item Often, other governmental institutions hold tax-related information but these data are not shared with the IRS.
    \item For example, millions of people receive health insurance coverage through Medicaid or other governmental programs. But this information often stays within the health agencies and is not shared with the IRS. %Not receiving these data makes it harder for the IRS to make sure that people receive the correct amounts of tax credits.
    \item Another example is that local governments keep track of who owns a house and how much property tax they pay. But they do not send this information to the IRS. %This means the IRS does not automatically know if someone owns a home or rents it out, which would be relevant for the IRS to assess housing-related tax deductions.
    \item For this data, the IRS would not be unable to prepopulate.  Therefore, taxpayers would have to declare it themselves.
    \item It may be more efficient if all tax-relevant information were consistently passed on to the IRS. This would help the IRS to offer prepopulated tax returns.
    \item The IRS could potentially save US taxpayers a fortune in tax filing costs by providing prepopulated tax returns. The US tax administration, the Internal Revenue Service, estimates that taxpayers need an average of 13 hours and \$270 for tax software and assistance to file their taxes.
    \item Moreover, most taxpayers would not lose much financial data privacy. Most tax-relevant data is already held by institutions of the government. In any case, honest taxpayers are already reporting all their tax-relevant data to the IRS. 
    \item There is an ongoing debate about whether the IRS should provide US taxpayers with free tax software and pre-populated tax returns. 
    \item But an organized group consisting of tax preparers and tax software providers is actively lobbying against this at the federal level. They are generating much revenue by selling tax software licences which they fear to loose when the government introduces their own tax filing services. 
   \item Until now, they have been able to prevent a major change to the tax filing system and secure their profits.
   \item For the individual taxpayer, however, a change to the tax filing system would certainly be beneficial.
\end{itemize}
%\begin{itemize}
%    \item Governmental institutions collect data on citizens in order to fulfill their functions and provide their services.
%    \item For instance, the Internal Revenue Service does not rely solely on taxpayers' self-reported information. It also collects information reported by third parties to verify the correctness of the self-reported information.
%    \item Third-party reporting means the tax authority uses information from other sources --- like employers, banks or insurance companies. This information is automatically sent from domestic firms and institutions to the tax authority.
%    \item For example, employers report to the tax administration how much they pay their employees in wages.
%    \item Furthermore, banks automatically inform the tax authority about interest and dividend payments of assets held by a person in a bank.
%    \item But the IRS does not automatically obtain all tax-relevant data it could use. 
%    \item Often, other governmental institutions hold tax-related information but these data are not shared with the IRS.
%    \item For example, millions of people receive health insurance coverage through Medicaid or other governmental programs. But this information often stays within the health agencies and is not shared with the IRS. Not receiving these data makes it harder for the IRS to make sure that people receive the correct amounts of tax credits.
%    \item Another example is that local governments keep track of who owns a house and how much property tax they pay. But they do not send this information to the IRS. This means the IRS does not automatically know if someone owns a home or rents it out, which would be relevant for the IRS to assess housing-related tax deductions.
%    %\item But this is tax-relevant information which is important for the IRS. 
%    \item For these data, the IRS needs to rely on taxpayers' declarations or to conduct tax audits verifying the information. This increases the likelihood of errors and leads to administrative costs and burden.
%    \item It may be more efficient if all tax-relevant information were consistently passed on to the IRS. This would help the IRS to fulfill its functions and provide services to taxpayers.
%    \item (For example, the tax authorities can better fight tax evasion if they have all tax relevant information.) 
%    \item (The tax authority would also be able to prepopulate tax returns if they have all necessary data. Then taxpayers would not need to fill in the tax returns themselves but only check and confirm the tax returns prepopulated by the tax authority.) 
%\end{itemize}

\subsubsection*{Information Text}
\subsection{Full Info Treatment}
\subsubsection*{General}
The last information treatment includes information from all treatments. It can be interpreted as a balanced video including the highlights of all the treatments. It helps us to understand how information transmission is changing if the information is presented next to each other. 
\subsubsection*{Information Text}
\begin{itemize}
    \item One problem governments face is the loss of tax revenue due to tax evasion. This means that some taxpayers deliberately hide some of their income or wealth from the tax administration. 
    \item In the US, for example, the Internal Revenue Service estimates that only 85\% of taxes are paid voluntarily and on time. The remaining 15\% amount to \$696 billion and more than \$600 billion of these are never collected
    \item Research has shown that this big amount of evaded taxes is mainly result of tax evasion among wealthy taxpayers.
    \item That means a small group of the population is responsible for a large part of the missing taxes. For example (x)\% of missing taxes are due to the (y)\% richest part of the population. 
    \item One tool to combat tax evasion is third-party reporting.
    \item Third-party reporting means the tax authority uses information from other sources --- like employers, banks or insurance companies. This information is automatically sent from domestic firms and institutions to the tax authority.
    \item For example, employers report to the tax administration how much they pay their employees in wages.
    \item Furthermore, banks automatically inform the tax authority about interest and dividend payments of assets held by a person in a bank. 
    \item Besides domestic third-party reporting, there is also the international exchange of information on bank data of foreign tax residents. %This is called automatic exchange of information and it has proven to be effective in reducing tax evasion.
    \item Increasing the amount of data that is third-party reported is an effective way to reduce tax evasion.
    \item Currently, the IRS does not automatically obtain all tax-relevant data it could use.
    \item Often other governmental institutions hold tax-related information but these data are not shared with the IRS.
    \item For example local governments keep track of who owns a house and how much property tax they pay. But they do not send this information to the IRS. This means the IRS does not automatically know if someone owns a home or rents it out, which would be relevant for the IRS to assess housing-related tax deductions.
    \item For these data, the IRS needs to rely on taxpayers’ declarations or to conduct tax audits verifying the information
    \item It may be more efficient if all tax-relevant information were consistently passed on to the IRS.
    \item Not only would this help the IRS to reduce tax evasion, it would also allow them to potentially prepopulate tax returns. 
    \item This means that the tax returns are pre-filled with information the tax administration already has. Prepopulated tax returns make the process of filing taxes much shorter and easier, with lower potential for making mistakes.
    \item The IRS could potentially save US taxpayers a fortune in tax filing costs by providing prepopulated tax returns. The US tax administration, the Internal Revenue Service, estimates that taxpayers need an average of 13 hours and \$270 for tax software and assistance to file their taxes.
    \item Moreover, most taxpayers would not lose much financial data privacy. Most tax-relevant data is already held by institutions of the government. In any case, honest taxpayers are already reporting all their tax-relevant data to the IRS. 
     \item (However, an organised group consisting of tax preparers and tax software providers is actively lobbying against a change to the government-provided tax filing services. Even though this change might benefit taxpayers, the group is protecting its big revenue from selling tax software licences and tax preparation services.)
\end{itemize}
\subsubsection*{Follow up questions}
We will ask 2-3 very basic questions to understand if respondents actually watched all parts of the video.

\newpage
\printbibliography

\end{document}